
\subsection{A környezet}
Az Egyetemen működő kutatócsoport a LoRaWan protokoll kiberbiztonsági kutatását végzi.

A tesztlabor egy, az Open Telekom Cloud-on futó virtuális gépen található. A kutatáshoz használt komponensek (LoRaWan Gateway, ELK-stack, saját service-ek) docker környezetben futnak.



\subsection{A cél}

A korábbi működés szerint a forrásokat a fejlesztő eljuttatja a szerverre, majd a Docker Compose indításakor a forrásból buildeli le.

A projektmunka által megvalósítani kívánt eredmény az, hogy a verziókezelő rendszerbe érkezett push hatására az új állapot automatikusan telepítésre kerüljön. A munka hatáskörébe ennek a rendszernek az alapvető működésének a kiépítése tartozik, nem a végleges, produkció-kész, finomhangolt állapot. Annak kialakítása a rendszer hosszabb távú használata során szerzett tapasztalatokra, és a kutatócsoport tagjaival való egyeztetések eredményeire épülhet.

A diplomamunkám során a feladatom lesz több új microservice készítése, valamint a meglévők átfogó keretrendszerbe foglalása. Ebben a folyamatban jelentős segítséget jelent egy megbízható CI/CD folyamat.


